\par\bigskip\pdfbookmark[0]{De nexu. Fragmentum VI}{nexus-fragment-VI}\centerline{[VI.]}\medskip
\textsc{Theorema.} Multitudo classium, in quas omnes formae binariae proprie primitivae determinantis negativi $-D$\begingroup\renewcommand{\thefootnote}{*)}\footnote{[distribuuntur.]}\endgroup, aequalis est
\[
\frac\pi4\times\surd D\times\text{Prod. ex. }\frac{p-1}p\frac{q+1}q\times\frac{rr-1}{rr}
\]
designantibus
\begin{description}
\item[$p$] omnes numeros primos\begingroup\renewcommand{\thefootnote}{**)}\footnote{[impares.]}\endgroup\ quorum non-res. est $-D$
\item[$q$] omnes numeros primos\textsuperscript{**)} quorum res. $-D$
\item[$r$] omnes numeros primos\textsuperscript{**)} ipsum $D$ metientes
\end{description}
\[
=\frac{\frac\pi4\surd D\ \text{Prod. ex }\frac{rr-1}{rr}}{1\pm\tfrac13\pm\tfrac15\text{ etc.}}
\]

ubi in denom. signum posit. praeponitur fractt., quarum denom. sunt in forma non divis.; negat. iis, quarum denom. sunt in forma divisorum ipsius $xx+D$; eae vero, quarum denom. ad $D$ non forent primi, omnino omittuntur\begingroup\renewcommand{\thefootnote}{***)}\footnote{[Bezeichnet man mit $m$ alle positiven ganzen Zahlen, die relative Primzahlen zu $2D$ sind, und benutzt man das durch \textsc{Jacobi} verallgemeinerte Symbol von \textsc{Legendre}, so ist die obige Regel für die Zeichenbestimmung in folgender Weise zu berichtigen: in der vorhergehenden Formel ist der Nenner
\[
1\pm\tfrac13\pm\tfrac15\text{ etc.}=\Sigma\pm\left(\frac{-D}m\right)\frac1m
\]
wo das obere oder untere Zeichen zu nehmen ist, je nachdem die Zahl $m$ ein Product aus einer geraden oder ungeraden Anzahl (gleicher oder ungleicher) Primzahlen ist; dagegen ist im Zähler der nachfolgenden Formel
\[
1\pm\tfrac13\pm\tfrac15\ldots=\Sigma\left(\frac{-D}m\right)\frac1m
\]
]}\endgroup.

\SourceLeaf{296}
\[
\begin{aligned}
&=\frac{2\surd D(1\pm\tfrac13\pm\tfrac15\ldots)}\pi\\
&=\frac{\operatorname{cotg}\theta\pm\operatorname{cotg}3\theta\pm\operatorname{cotg}5\theta\ldots\pm\operatorname{cotg}n\theta}{N:\surd D}
\end{aligned}
\]
ponendo $\theta=\frac\pi N$, $N=\left\{\begin{smallmatrix}1\\4\end{smallmatrix}\right\}D$ et ponendo pro $n$ omnes numeros ad $D$ primos signo ut supra determinato\begingroup\renewcommand{\thefootnote}{*)}\footnote{[Siehe die weiter unten folgende Note zu diesem Fragment.]}\endgroup.

Pro determ. pos. erit mult. Classium\begingroup\renewcommand{\thefootnote}{**)}\footnote{[In der nachfolgenden Formel bedeutet $D$ den positiven Determinanten, und es ist
\[
1\pm\tfrac13\pm\tfrac15\ldots=\Sigma\left(\frac Dm\right)\frac1m
\]
]}\endgroup
\[
=\frac{2\surd D(1\pm\tfrac13\pm\tfrac15\ldots)}{\log .T+U\surd D}
\]
Designantibus $T,U$ valores minimos quantitatum $t,u$ aequationi $tt-Duu=1$ satisfacientes
\[
=\frac{\log\sin\tfrac12\theta\pm\log\sin\tfrac32\theta\pm\log\sin\tfrac52\theta\text{ etc.}}{\log T+U\surd D}
\]
\BookRule{.30\textwidth}
